% Empirical Strategy

\subsection{Empirical strategy} \label{subsec:empirical-strategy}

We examine how short selling demand signals in the equity lending market affect mutual fund portfolio allocation. Consider the following general characterization of mutual fund holdings, where $Y_{i,j,t}$ is the holding for fund $j$ in stock $i$ at time $t$:
\begin{align} \label{eq:equil_allocation}
Y_{i,j,t} & = f\left(Y_{i,j},\; Y_{i,t},\; Y_{j,t},\; \phi_{i,j,t} \right)
\end{align}
Mutual fund holdings are an equilibrium outcome determined by four components: a fund-stock time-invariant component $Y_{i,j}$, which reflects a fund's persistent preferences for a stock based on investment philosophy or style; a common stock-time component $Y_{i,t}$, which captures signals common to all investors including aggregate short selling demand; a fund-time component $Y_{j,t}$, which captures fund-level investment style, portfolio breadth, and governance; and a fund-stock-time specific component $\phi_{i,j,t}$.

Our focus is on the fund-stock-time component $\phi_{i,j,t}$: how fund $j$'s specific demand for stock $i$ changes when the fund receives a \textit{private} short selling demand signal $OnLoan_{i,j,t-1}$ through the equity lending market. For example, consider a fund active in the lending market that lends $OnLoan_{i,j,t-1}$ shares of stock $i$ at time $t-1$. The amount on loan reveals information about short selling demand that is not contained in public signals---because public signals are absorbed by $Y_{i,t}$---but is observed privately by the fund through its lending contracts.\footnote{A measure of aggregate short selling demand is available to institutional investors. Brokers often provide data on special lists and borrowing costs to institutional investors with substantial trading volume. Markit reciprocally shares lending market information with participating clients.}

In a saturated fixed effect estimation, we estimate changes in fund-stock portfolio holdings as a function of a fund's lending activity, including fund-stock, stock-time, and fund-time fixed effects ($\alpha_{i,j}$, $\alpha_{i,t}$, and $\alpha_{j,t}$). These fixed effects absorb public stock-time price signals common to all investors, fund-time information and investment decisions common to all stocks in a fund's portfolio, and fund-stock time-invariant preferences.

To identify the effect of the private signal $OnLoan_{i,j,t-1}$ on portfolio allocation, we estimate:
\begin{align} \label{eq:baseline}
Y_{i,j,t} & = \alpha_{i,j} + \alpha_{i,t} + \alpha_{j,t} + \beta_{1}\, \mathit{OnLoan}_{i,j,t-1} + \mathbf{X}_{i,j,t-1}'\,\gamma_1 + \epsilon_{i,j,t},
\end{align}

\noindent where $Y_{i,j,t}$ is the monthly percentage change in the number of shares of stock $i$ held by fund $j$ at time $t$ ($\mathit{PosChange}_{i,j,t}$). $\mathit{OnLoan}_{i,j,t-1}$ is the dollar value of shares lent by fund $j$ in stock $i$ at $t-1$ divided by the market capitalization of stock $i$, standardized to mean zero and unit standard deviation to facilitate comparison of effect magnitudes across specifications.\footnote{Standardization is performed within the estimation sample. All $\mathit{OnLoan}$ variables throughout the paper are standardized in this way unless stated otherwise.} This measures the fund's exposure to private short selling demand. $\mathbf{X}_{i,j,t-1}$ includes lagged position weight ($\mathit{PosWeight}_{i,j,t-1}$) and an Amihud~\citeyearpar{amihud:02} illiquidity rank ($\mathit{ILLIQ}_{i,j,t-1}$). Standard errors are clustered at the fund level.\footnote{For the baseline specification in equation~(\ref{eq:baseline}), clustering at the fund level is appropriate because the primary source of residual dependence is within-fund correlation across the security and time dimensions; the identifying variation is fund-specific rather than time-specific. For specifications involving the mandatory disclosure events in Sections~\ref{sec:information_channel}--\ref{sec:market_structure}, we double-cluster by fund and time. Disclosure events are correlated across funds in the same month, introducing a time-series dimension to the error structure that a single fund cluster does not fully capture.} Variable definitions are provided in Appendix~\ref{app:variables}.

The coefficient $\beta_1$ identifies the effect of the private signal: we predict $\beta_1 < 0$ if lenders act on the information embedded in short selling demand. The counterfactual is funds that hold the same stock but do not participate in lending and therefore receive no private signal. Any effect of aggregate short selling demand $OnLoan_{i,t}$ is absorbed by $\alpha_{i,t}$.

\subsection{Active vs.\ passive funds} \label{subsec:active-passive}

A critical identification test is whether the trading response to lending signals is confined to active fund managers who can freely adjust portfolios, rather than passive fund managers who must track a benchmark index. We estimate:
\begin{align} \label{eq:baseline_active}
Y_{i,j,t} & = \alpha_{i,j} + \alpha_{i,t} + \alpha_{j,t} + \beta_{1}\, \mathit{OnLoan}_{i,j,t-1} + \beta_{2}\, \mathit{OnLoan}_{i,j,t-1} \times \mathit{Active}_{j} + \mathbf{X}_{i,j,t-1}'\,\gamma_1 + \epsilon_{i,j,t},
\end{align}

\noindent where $\mathit{Active}_j$ equals one for actively managed funds. The coefficient $\beta_1$ captures the response for passive funds, which we expect to be zero since passive funds cannot trade on information. $\beta_2$ captures the incremental response for active funds relative to active non-lenders; our information acquisition hypothesis predicts $\beta_2 < 0$. Note that aggregate short selling demand $\mathit{OnLoan}_{i,t-1}$ is absorbed by the stock$\times$time fixed effect $\alpha_{i,t}$, so the coefficients identify the fund-specific private signal over and above any aggregate demand common to all funds.

\subsection{Identification via mandatory short position disclosures} \label{subsec:identification-disclosure}

Our primary identification strategy exploits mandatory public disclosures of large short positions. Under Article~9 of Regulation (EU) No~236/2012, any net short position exceeding 0.5\% of issued share capital must be publicly disclosed the following trading day. Above 0.5\%, further disclosures are required at 0.1\% increments. The regulation applies to all investors irrespective of origin, with standardized disclosure across EU countries.\footnote{The regulation and relevant data are described in detail in \citet{jank_smajlbegovic:2015}, \citet{jones_reed_waller:16}, and \citet{jank_roling_smajlbegovic2020}. Publicly disclosed short positions are more frequently observable for large and liquid stocks \citep{jank_roling_smajlbegovic2020}, consistent with our sample.}

We focus on \textit{new} disclosure events---stocks that have not had any public short position disclosure in the preceding six months. This restriction isolates the most informative public signals: follow-on disclosures for stocks with known short positions are less informative \citep{jones_reed_waller:16}. These disclosed positions predict negative returns \citep{jank_smajlbegovic:2015} but do not fully reveal the private structure of shorting demand below the disclosure threshold \citep{jank_roling_smajlbegovic2020}.

\textbf{Selection concerns.} The main identification challenge is that lending activity may be endogenous to portfolio decisions. A fund might reduce a position precisely because it can no longer lend shares (liquidity-driven), or might lend more shares in positions it plans to reduce (selection). The three-way fixed effects structure directly addresses this concern: $\alpha_{i,t}$ absorbs stock-time common shocks including aggregate borrowing demand, $\alpha_{j,t}$ absorbs fund-level portfolio rebalancing, and $\alpha_{i,j}$ absorbs time-invariant fund-stock preferences. The identifying variation is the idiosyncratic deviation of a fund's lending in a specific stock at a specific time, relative to what is predicted by these three layers of fixed effects.

A related mechanical concern is that the fund$\times$security fixed effect $\alpha_{i,j}$ demeans $OnLoan_{i,j,t-1}$ within each cell, and because lending is zero in most observations (only 2\% of fund-security-month observations have any shares on loan), the demeaned residual is positive in lending months and negative otherwise. If a fund mechanically reduces positions as lending terminates---with no informational content---this autocorrelation structure could generate a spurious negative $\hat{\beta}_1$. Three features of the design collectively rule out this mechanical story. First, and most directly, passive funds are subject to the identical fund$\times$security fixed effect structure and have substantially \emph{higher} lending rates (69.8\% of passive funds participate in lending versus 25.5\% of active funds), yet they exhibit precisely zero trading response (Table~\ref{table_disclosure}, column~8). A mechanical channel would affect passive funds at least as strongly as active funds, since the demeaning algebra is identical and passive funds' lending is more prevalent. Second, the disclosure timing interaction is forward-looking: the mechanical channel predicts no amplification when lending in month $t-1$ precedes a \emph{future} disclosure event in month $t+1$, because the mechanical story has no access to forward-looking public information. Yet the interaction $OnLoan_{i,j,t-1} \times \mathit{Disclosure}_{i,t+1}$ is the strongest and most significant term in Table~\ref{table_timing}. Third, the mechanical story predicts symmetric effects around any systematic shock to the stock---not the asymmetric lead-only pattern documented in Section~\ref{sec:information_channel}. We plan to provide formal diagnostics (excluding lending-termination observations; AR(1) residual tests) in the revision.

We further address endogeneity by exploiting the timing of mandatory disclosures. If the lending signal is informative, active lenders should reduce positions \textit{before} a large short position is publicly disclosed---not after. We test this using the following regression:
\begin{align} \label{eq:big-shorts}
Y_{i,j,t} & = \beta_1\, \mathit{OnLoan}_{i,j,t-1} + \beta_2\, \mathit{Disclosure}_{i,t+1} \times \mathit{OnLoan}_{i,j,t-1} \nonumber\\
& \quad + \beta_3\, \mathit{Disclosure}_{i,t} \times \mathit{OnLoan}_{i,j,t-1} + \beta_4\, \mathit{Disclosure}_{i,t-1} \times \mathit{OnLoan}_{i,j,t-1} \nonumber\\
& \quad + \mathbf{X}_{i,j,t-1}'\,\gamma + \alpha_{i,t} + \alpha_{j,t} + \alpha_{i,j} + \epsilon_{i,j,t},
\end{align}
where $\mathit{Disclosure}_{i,t}$ is an indicator equal to one if there is a \textit{new} public short position disclosure in month $t$ for stock $i$. $\mathit{Disclosure}_{i,t+1}$ and $\mathit{Disclosure}_{i,t-1}$ are the lead and lag of this indicator, respectively. The key coefficient is $\beta_2$: lenders reduce positions in the month \textit{before} disclosure. We expect $\beta_3 = 0$ and $\beta_4 = 0$ because lenders' information advantage evaporates once positions are publicly disclosed.

This timing pattern---pre-announcement selling by lenders but not by non-lenders---is inconsistent with mechanical explanations (e.g., recall of lent shares to vote) and is most consistent with active information acquisition through the lending process. Standard errors are double-clustered at the fund and time level for all disclosure-sample specifications.

\textbf{Contribution.} Our approach builds on and extends the existing literature. \citet{EvansFerreiraPrado2017} document that securities lending generates revenues for mutual funds but do not examine whether lending conveys information. The theoretical framework in \citet{Nurisso2026} provides our primary identification motivation: active lenders acquire private information but face a commitment problem, so that the information value of lending should be detectable precisely in the window before private information becomes public, and should be larger in concentrated markets where each lender's signal is more precise. We provide empirical confirmation of these theoretical predictions using three identification advances. First, we use European regulatory data to directly link individual lending contracts to portfolio decisions, isolating the fund-level private signal. Second, we use mandatory disclosure events as the key external validation: if active lenders exploit private information, position reductions should be concentrated immediately \emph{before} the information becomes public---the direct test of the Nurisso information channel. Third, by examining market structure---lending concentration, active agent counts, and portfolio breadth---we identify conditions under which the information value of the lending signal is highest, consistent with Nurisso's result that competition among lenders governs information precision. These market structure results are novel to the literature.
